% !PRINTED 317
\weberpage{317}
%<<CONT
fach, besonders in der Tschirnhausen-Transformation gute Dienste
geleistet haben.

Wir haben diese Functionen im \S.\,4 definirt durch
\[
\begin{aligned}
f_0(x) &= 1\\
f_1(x) &= x+a_1\\
f_2(x) &= x^2+a_1x+a_2\\
&\ \cdots\cdots\cdots\cdots\cdots\cdots\cdots\cdots\\
f_{n-1}(x) &= x^{n-1}+a_1x^{n-2}+a_2x^{n-3}+\cdots+a_{n-1}\\
f(x)=f_n(x) &= x^n+a_1x^{n-1}+a_2x^{n-2}+\cdots+a_{n-1}x+a_n.
\end{aligned}
\]
und dafür die Recursionsformel aufgestellt
\[
f_v(x)-xf_{v-1}(x)=a_v.
\]
Man hat nun, wie schon an der angeführten Stelle gezeigt ist,
\weq{3}{f(x)=(x-\alpha)\left[x^{n-1}f_0(\alpha)+x^{n-2}f_1(\alpha)+\cdots+f_{n-1}(\alpha)\right]+f(\alpha).}
Der Anblick dieser Formel zeigt, dass, wenn $\alpha$ eine Zahl ist, die
die Functionen
\weq{4}{f_0(\alpha),\ f_1(\alpha),\ f_2(\alpha)\ldots f_{n-1}(\alpha),\ f(\alpha)}
positiv macht, kein positives $x$ grösser als $\alpha$ existiren kann,
was $f(x)$ zu Null macht, da dann auf der rechten Seite von (3) lauter
positive Glieder stehen.

Wenn also ein positives $\alpha$ so bestimmt ist, dass die Functionen
(4) alle positiv sind, so ist dies eine obere Grenze für die positiven
Wurzeln von $f(x)$.

Ebenso erhält man eine untere Grenze $-\beta$ für die negativen Wurzeln,
wenn man die positive Grösse $\beta$ so bestimmt, dass
\weq{5}{f_0(-\beta),\ -f_1(-\beta),\ f_2(-\beta),\ \ldots (-1)^n f(-\beta)}
alle positiv werden.

Diese Bestimmung der Grenzen ist zwar in der Rechnung einfacher zu
handhaben, giebt aber doch unter Umständen minder genaue Grenzen als
die \gesperrt{Newton}'sche Bestimmung.

Wenn man nämlich (3) fortgesetzt nach $x$ differentiirt und dann
$x=\alpha$ setzt, so erhält man unter der Voraussetzung, dass die
Grössen (4) alle positiv sind, für die Derivirten $f'(x)f''(x)\ldots$
Ausdrücke, die für $x=\alpha$ aus lauter positiven Gliedern bestehen.
Wenn also die Ausdrücke (4) positiv sind, so sind nothwendiger Weise
auch die Derivirten (1) alle positiv, während das Umgekehrte nicht
nothwendig der Fall ist.
